\begin{textsample}{POS Dim 8 – global\_north\_2025\_03 – Score 13.00 – global\_north\_2025\_03/122519943.txt}  \label{ex:f8_pos_084}
The Labour Muslim Network, an organisation which promotes British Muslim engagement with the Labour Party, has slammed Prime Minister Keir Starmer for rowing back on Foreign Secretary David Lammy ' s statement on Monday that Israel was breaking international humanitarian law.
Asked about the blockade and its consequences in parliament on Monday, Lammy said :"This is a breach of international law.
Israel quite rightly must defend its own security.
But we find the lack of aid -- it ' s now been 15 days since aid got into Gaza -- unacceptable, hugely alarming and very worrying."But Starmer ' s official spokesperson said on Tuesday afternoon that the UK government position remains that Israel ' s actions in Gaza are"at clear risk"of a breach.
The Labour Muslim Network ( LMN ) told Middle East Eye the move was"shocking"and"disgraceful"."The horrific and barbaric attacks by Israel on innocent civilians overnight in Gaza are yet another flagrant disregard of international and humanitarian law,"a @ @ @ @ @ @ @ @ @ @ : Jerusalem Dispatch Sign up to get the latest \textbf{insights} and analysis on Israel-Palestine, alongside Turkey \textbf{Unpacked} and other \textbf{MEE} \textbf{newsletters}"No British government with a conscience should continue to call Israel an ally or continue any form of military or economic relationship with an Israeli government that consistently and willfully disregards international law."They added :"Foreign Secretary David Lammy was right to call the blockade of goods a breach of international law in the House of Commons.
Attempts to row back and retract this assertion since are as shocking as they are disgraceful."' No British government with a conscience should continue to call Israel an ally ' - Labour Muslim Network The statement comes after Israel killed more than 400 Palestinians without warning on Tuesday morning, unilaterally ending its ceasefire with Hamas.
Tuesday ' s attack, which targeted Palestinians displaced by the conflict, was one of the bloodiest bombardments of the Gaza Strip since the war began.
The Labour government had not accused Israel of breaking international @ @ @ @ @ @ @ @ @ @ in parliament on Monday represented a significant shift in the UK ' s position. ' Deliberate obfuscation ' Outrage has been mounting in parliament at Starmer ' s move.
Independent \textbf{MP} Ayoub Khan told \textbf{MEE} that Number 10 ' s position was"shameful", accusing the government of"deliberate obfuscation"."This pattern of rhetoric and failing to condemn actions while refusing to impose any meaningful consequences has emboldened the Israeli government to continue its campaign of murder, destruction, starvation, and \textbf{displacement} in Gaza,"he said."The government would do well to align itself with the new found moral compass of its Foreign Secretary rather than undermine him."The Council for Arab-British Understanding said :"In a disgraceful move, No 10 Downing Street has tried to retract the Foreign Secretary ' s clear statement accusing Israel of violating international law with its blockade of Gaza."It undermines the government ' s claims that it respects international law and exposes a clear chasm between the Prime Minister and his Foreign @ @ @ @ @ @ @ @ @ @ the government ' s move as"nothing short of appalling"."The UK Government must condemn these crimes in the strongest terms immediately and stop its complicity in this catastrophic crisis."Middle East Eye delivers independent and \textbf{unrivalled} \textbf{coverage} and analysis of the Middle East, North Africa and beyond.
To learn more about \textbf{republishing} this content and the \textbf{associated} \textbf{fees}, \textbf{please} \textbf{fill} out this form.
More about \textbf{MEE} can be found here.

% matched lemmas: associate, coverage, displacement, fee, fill, insight, mee, mp, newsletter, please, republish, unpack, unrivalled
\end{textsample}
